\documentclass[13pt]{ctexbeamer}


\usepackage{amsmath,amssymb,amsthm}             % AMS Math
\usepackage{graphicx}
\usepackage{epstopdf}
\usepackage{tikz}
\linespread{1.3}

\usepackage{mathrsfs}  %花写字母


%%%=== theme ===%%%
%\usetheme{Warsaw}
%\usetheme{Copenhagen}
%\usetheme{Singapore}
\usetheme{Madrid}
%\usefonttheme{professionalfonts}
%\usefonttheme{serif}
% \usefonttheme{structureitalicserif}
%%\useinnertheme{rounded}
%%\useinnertheme{inmargin}
\useinnertheme{circles}
%\useoutertheme{miniframes}
\setbeamertemplate{navigation symbols}{}
%\setbeamertemplate{footline}[page number]
\setbeamertemplate{footline}[frame number]


\titlegraphic{\includegraphics[width=2cm]{tjnu.jpg}}


\usepackage{color}
\definecolor{linkcol}{rgb}{0,0,0.4}
\definecolor{citecol}{rgb}{0.5,0,0}

\usepackage{xcolor}
\newcommand{\red}[1]{\textcolor{red}{#1}}
\newcommand{\blue}[1]{\textcolor{blue}{#1}}
\newcommand{\green}[1]{\textcolor{green}{#1}}


  \usepackage{graphicx}
  \DeclareGraphicsExtensions{.eps}
%   \usepackage[a4paper,pagebackref,hyperindex=true,pdfnewwindow=true]{hyperref}


\begin{document}



\title[]{论文写作指导 --- 数学文章的结构}
\author[]{{\large 张彪} }
\institute[]{{\normalsize
		天津师范大学\\[6pt]
		zhang@tjnu.edu.cn}}

\date{}



\AtBeginSection[]
{
\begin{frame}
	\frametitle{数学文章的结构}
	\tableofcontents[currentsection]
    \setcounter{tocdepth}{1}
\end{frame}
}



\begin{frame}
\maketitle
\end{frame}



\begin{frame}{第一章~数学文章的结构}
一篇数学论文一般由如下几个部分组成：
%\begin{itemize}
%\item {题目}
%\item 摘要
%\item 引言
%\item 主体
%\item 结论
%\item 致谢
%\item 文献
%\item 附录
%\end{itemize}
\setcounter{tocdepth}{1}
	\tableofcontents
\end{frame}


\section{题目}
\begin{frame}{题目}
题目是文章的一句广告语。
\begin{itemize}
	\item 题目不宜太长，也不宜太短。
	\item 题目应具体、明确，反映文章的主要贡献。
	\item 不要泛指一个过大的方向。
\end{itemize}
\end{frame}


\begin{frame}{题目}
	题目是文章的一句广告语。
	\begin{itemize}
		\item 题目不宜太长，也不宜太短。
		\item 题目应具体、明确，反映文章的主要贡献。
		\item 不要泛指一个过大的方向。
	\end{itemize}

题目是文章的脸面。文章的题目需：
\begin{itemize}
	\item 含有一定的信息量，使读者课通过搜索引擎得到你的文章；
	\item 能够吸引读者的注意力；
	\item  言简意赅地表述文章的研究成果；
	\item 和现有的标题有一定的区别。
\end{itemize}
\end{frame}

\section{摘要}
\begin{frame}{摘要}
摘要是一篇有足够信息的微型文章。	文章的摘要应该：
	\begin{itemize}
		\item 概括文章的主要目的、思想和结果。
		\item 极可能简明扼要，但又要有足够的内容。
		\item 用词精确、意思明确，尽可能让更多的人读懂你的叙述。
	\end{itemize}

摘要是文章的心脏。它一般要回答
	\begin{itemize}
		\item 本文要干什么？本文涉及的问题是什么？
		\item 研究的问题如何解决？研究方法是什么？
		\item 主要的结果是什么？问题彻底解决还是部分解决了？
		\item 研究结果的意义？对科学或对读者有多大帮助？
	\end{itemize}

\end{frame}

\section{关键词和学科分类}
\begin{frame}{关键词和学科分类}
	关键词keywords
	\begin{itemize}

		\item 关键词的设置是为了让数学论著归档服务机构（如美国《数学评论》）能把你的论文放到其门类齐全的数据库中的正确位置。

		\item 科学引用指标检索工具（SCI）利用关键词给文章分门别类，而读者根据关键词检索与它研究或兴趣相关的文献。
		\item 最好只写单数形式，但是写成复数也无伤大雅
	\end{itemize}
\end{frame}


\begin{frame}{关键词和学科分类}
	数学学科分类 Mathematics Subject Classification  （MSC)
	\begin{itemize}

		\item 分类系统是由美国《数学评论》Math Reviews（MR）和德国《数学文摘》（zbMath）共同制定的。2010，2020

		\item https://mathscinet.ams.org/mathscinet/msc
	\end{itemize}
\end{frame}


\section{引言}
\begin{frame}{引言}
	引言应该比较全面、准确并且客观地介绍文章中将要讨论
	\begin{itemize}
		\item 问题的背景材料和简要发展，
		\item 人们对此已做的相关贡献，
		\item 写这篇文章的动机，
		\item 本文的主要结果。
	\end{itemize}
需要注意：
	\begin{itemize}
	\item 引言的客观性非常重要，切记自我吹嘘。\\
	~~~~~~~~~~~~``文章在自己~~评价在别人''~~~~(华罗庚)
	\item 不要随意贬低别人的研究。
	\item 不要把别人的甚至自己以前的文章句子原封不动地直接移植过来。
\end{itemize}
\end{frame}

\begin{frame}{1~引言的开首}
		\begin{itemize}
		\item ``好的开始是成功的一半'

		\item 最糟糕的开头可能是先给出一堆数学符号和定义。

		\item  最好的引言应该这样开头：
	\begin{itemize}
		\item 先做一些较为通俗浅显的描述，这样会给读者一个容易的起头，以便他们轻松自如、有条不紊地登堂入室，进入角色，并能提高继续读下去的兴趣。
		\item  把你所关心的问题提出来，开门见山、直奔主题，不要不着边际地叙述与文章的中心关系不大的事物。要用直截了当、浅显易懂的语音把问题指出来。
	\end{itemize}
	\end{itemize}
\end{frame}


\begin{frame}{2~引言的中间}
	\begin{itemize}

		\item 定义问题

		\item 解释准备解决什么问题
		\item 总结以前取得的主要成果及不足之处
		\item 简介一下你解决问题主要的手法
		\item 指出写作的动机和目的是必要的。
	\end{itemize}
\end{frame}

\begin{frame}{3~引言的结尾}
	\begin{itemize}

		\item 简单叙述文章的组成部分

		\item 对于文章后面的每一节内容写上一句话，概括地告诉读者这一节是干什么的
		\item 用不同的语句给出每一节大意一个简单的小结。
	\end{itemize}
\end{frame}

\section{主体}

\begin{frame}{主体}




	Preliminaries（预备知识）
	\begin{itemize}
		\item  引进符号说明和概念定义。
		\item  有些已知的但对这篇文章有用的结果也需要单独引出来。这样做是为下一节作准备，并能区别什么是你文章的主要新结论，什么是以前的结果。
	\end{itemize}
	Main results（主要结果）
		\begin{itemize}
		\item  1-2个主要定理，最多3-4个
		\item 定理的叙述中要包含完整的条件和结论
	\end{itemize}

Proofs（证明）
\end{frame}


	\setcounter{tocdepth}{2}
\frame{\tableofcontents[currentsection]}
%\frame{\tableofcontents[sectionstyle=shaded , currentsection]}

\subsection{文章的章节标题}
\begin{frame}{文章的章节标题}
章节的标题是文章的骨骼。

对读者来说，文章的结构或小标题

\begin{itemize}
    \item  使读者更容易找到他们感兴趣的部分;
\item 使读者更容易找到作者的主要贡献或文章的亮点;
\item 使读者从逻辑上更易理解作者的写作意图。
\end{itemize}


而对作者来说, 章节小标题
\begin{itemize}
    \item  对相应的章节起到画龙点睛的作用;

\item 为自己写作前提供好的提纲;

\item 把文章划分成相对独立的部分使目的性强的读者直奔所需。
\end{itemize}
\end{frame}




\subsection{定义}
\begin{frame}{定义}
	\begin{itemize}
		\item
		制定一个定义时要考虑三个问题

		``\blue{为什么？}''

		``\blue{放哪里？}''

		``\blue{怎么样？}''
		\item
		首先，自问为什么你要做一个定义： 它是必须的么？
		\item
		不恰当地定义会使表达复杂化并且太多定义会压垮读者，因此设想自己正在为每一个定理支出一大笔花销是明智的。
		\item
		在给定的学科领域中符号是标准的，需要判断定义是否要给出。
		\item
		用多余的单词可以避免潜在的混淆。
	\end{itemize}
\end{frame}


\begin{frame}
	\begin{itemize}
		\item  第二个问题是``\blue{放哪里？}''
		\item 不推荐在一篇文章刚开始的地方放置一长串的定义。
		\item  一个定义应该被放置在该属于首次使用的位置。
		\item 如果定义较早给出，读者将不得不往回查找，这可能失去集中力（或者更糟糕，失去兴趣）。
		\item 尽量缩短定义和它首次被使用位置之间的距离。
	\end{itemize}
\end{frame}


\begin{frame}
	\begin{itemize}
		\item  为了强化几页没有使用的数学符号，你可以重复定义。

		\item 例如，``最佳步长$\alpha^{*}$如下''。

		\item  这个含蓄的再定义或者提醒了读者$\alpha^{*}$是什么, 或者再次确保他们已经准确地记住了它。
	\end{itemize}
\end{frame}


\begin{frame}
	\begin{itemize}
		\item  最后，一个术语如何被定义？
		\item 有可能只有唯一的定义或者有多种可能。
		\item 你应该以一个简短地、用基本性质或者基本思想的术语描述的，并且与相关定义相容的定义为目标。
		\item 举个例子，正规矩阵的标准定义为“矩阵$A\in \mathbb{C}^{n\times n}$ 满足$AA^{*}= A^{*}A$”。 至少有70种不同的方式定义正规性，但是没有一种比$AA^{*}= A^{*}A$简单容易。
	\end{itemize}
\end{frame}


\begin{frame}
	\begin{itemize}
		\item 依照惯例，在定义中if 的意思是 if and only if , 所以不要写

		``The graph is connected  \alert{if and only if }~there is a path from every node in $G$ to every other node in $G$''.

		而写成

		``The graph is connected \alert{if}  there is a path from every node in $G$ to every other node in $G$''.

	\end{itemize}
\end{frame}

\begin{frame}
	\begin{itemize}
		\item 把要定义的文字排版成斜体是惯例：

		``The graph is \emph{connected} {if}  there is a path from every node in $G$ to every other node in $G$~''.

		\item 这种强调也可以写为

		``The graph is defined to be {connected} {if}  \dots''

		或者

		``The graph is said to be {connected} {if}  \dots''
	\end{itemize}

\end{frame}


\subsection{例子}

\begin{frame}{例子}
	\begin{itemize}
		\item
		适应于技术写作（从教学到研究）所有形式的教学策略是在讨论一般情况前都会先讨论\blue{特殊例子}。
		\item
		尤其对于数学家来说，采取反证法是有吸引力的，但是\blue{以例子开头来解释说明是更有效的方法}。
%		\item
		能说明如何\blue{以特定实例开头}的一个很好的例子是Strang的《应用数学导论》第一章内容：

		{\small
			应用数学中最简单的模型是线性方程组，也是目前为止最重要的模型，我们以一个极为简单的例子开始本书的内容：
			\begin{align*}
				2 x_1+ 4 x_2& =2,\\
				4 x_1 +11x_2 & =1.
			\end{align*}
			在一些进一步的开场白之后, Strang继续详细地研究这个$2\times 2$矩阵和一个特殊的$4\times 4$矩阵。仅在数页之后就给出了一般的$n \times n$矩阵。}

		\item 课本上的\alert{练习题}是例子的一种形式。
	\end{itemize}


\end{frame}




\subsection{定理}
\begin{frame}{定理是什么？}


\begin{itemize}
	\item 定理
	\item 引理
	\item 命题
	\item 猜想
	\item 假设
\end{itemize}
\end{frame}

\begin{frame}
\begin{itemize}
\item \blue{定理}、\blue{引理}、\blue{命题}之间的差别是什么？
\item 在某种程度上，答案取决于结论在行文中的位置。
\item
一般地，\blue{定理}是一个具有独立意义的重要结论。
定理的证明通常是非平凡的。


\end{itemize}
\end{frame}




\begin{frame}{引理}

\begin{itemize}
	\item
	\blue{引理}是一个辅助的结果—是迈向定理的一个跳板。
	其证明可能容易也可能困难。

\item 一个明确的独立的值得概括但是不值得冠以定理头衔的结论也可被称为\blue{引理}。
事实上，有很多著名的引理。
\item
	一个结论是否应该被正式地规定为引理或者只是在文中简单地提及取决于你写作的\red{等级}。
\item
\alert{把所有结论都标注为定理是不可取的}，因为这样做就无法突出你文章的逻辑结构，
而且也无法将读者的关注点指向最重要的结论。
\item 如果你对一个结论是称为引理还是定理有所迟疑，那就称为\blue{引理}。
\end{itemize}


在一个线性代数的研究论文中，给出一个叙述``对称正定矩阵的特征值是正的'' 的引理是不恰当的，
因为这个权威的结论是众所周知的。

\end{frame}


\begin{frame}{命题}
\begin{itemize}
	\item
\blue{命题}比引理和定理使用范围要小并且他的意义也不尚明确。
它倾向于表示一个次要定理的方法。
\item 讲义和教科书的作者可能认为他温文尔雅的名字是的它比定理显得不那么令人生畏。
 \item 然而，命题并不是学生认为的“一个可能不正确的定理”。
\end{itemize}
\end{frame}

\begin{frame}{推论}
\begin{itemize}
\item \blue{推论}是引理、定理或者命题的直接或很容易推出的结论。
\item 区分推论和一个结论的延伸或者概括是非常重要的。
\item 当心不要过度颂扬一个推论，如给它错误的标注（定理），因为这将是它不恰当地得以突出而且会混淆主要结论（主定理）的地位。
\end{itemize}
\end{frame}

\begin{frame}
\begin{itemize}
\item 有多少结论被正式陈述为引理、定理、命题和推论是个人风格的问题。
\item 一些作者用一系列赋有定义以及评注的结论和定理来阐述他们的观点。
\item 另一种极端是，一些作者会极少正式地陈述结论。
\end{itemize}
\end{frame}

\begin{frame}{猜想}
\begin{itemize}
\item 第五种数学写作声明是\blue{猜想}。
\item 作者认为可能正确但是未曾被证明或反驳的陈述。
\item 作者通常会为了证实陈述的真实性写出一些强有力的证据。
\item  一个著名猜想的例子就是哥德巴赫猜想(1742年)，每一个大于2的偶数都可写成两个质数之和，它仍未被证明。
\item 提出一个猜想之后再反驳不一定是个坏事：确定这个猜想旨在回答的问题可能会是一个重要贡献。


\item 一个假设是进一步猜想的基础。
\item 站在自己立场上的假设是罕见的，如黎曼假设和连续统假设。
\end{itemize}
\end{frame}


\subsection{证明}

\begin{frame}{证明}
	\begin{itemize}
		\item 读者更想要知道大纲和关键的想法。
		\item 读者更希望学习一种可以应用在其他情况下的技术或原理。
		\item 当读者要详细研究证明的细节时，他们自然想花费最少的精力去了解它。
		\item 为了在这两种情况下帮助读者，强调一下
		\begin{itemize}
		\item \blue{证明的结构}
		\item \blue{每一步的难易}
		\item \blue{使证明成立的关键想法}
		\end{itemize}
		很重要。

	\end{itemize}
\end{frame}




\section{结论}
\begin{frame}{结论}

	\begin{itemize}
		\item  非常简短地叙述文章的主要贡献
		\item  解释你在引言里提到的问题和疑问。经过整篇文章的论证，也许可以谈论部分答案了
		\item 指出由于某些原因，文章没有考虑的其他方面或更广泛的问题,并说明本文的论证方法是否可以推广到这些情形
		\item  展望下一步，后续研究可以做什么？
		\item 如果合适，讨论一些和本研究相关的猜想
	\end{itemize}
\end{frame}

\section{致谢}
\begin{frame}{致谢}

	\begin{itemize}
		\item  致谢的对象是主要包括给文章写作提供过意见或者帮助的人或机构
		\item  很多研究基金规定获资助者在所发表的\alert{相关}论文中必须清楚表明获其资助，有时需列出所获资基金的号码。
		\item 不必提及匿名审稿人。但如果审稿人对改进文章的质量有很大的共享，应向他们致谢。
	\end{itemize}
\end{frame}

\section{参考文献}
\begin{frame}{参考文献}

	\begin{itemize}
		\item  文献引用
		\item  文献格式
	\end{itemize}
\end{frame}

\section{附录}
\begin{frame}{附录}

	\begin{itemize}
		\item  附录记载的主要是作者不想放在正文内的冗长复杂的定理证明。
		\item  另一种附录防止的是正文定理证明中所要的的一些预备证明或一些标准引理，甚至一些有关概念的定义等。
		\item 附录可以收入任何作者不便或不拟放在文章不要部分、但有不想随意丢掉的东西。
	\end{itemize}
\end{frame}




\end{document}